\documentclass[11pt]{article}

\usepackage[a4paper,margin=1in]{geometry}
\usepackage{amsmath}
\usepackage{amssymb}
\usepackage{booktabs}
\usepackage{enumitem}
\usepackage[hidelinks]{hyperref}
\usepackage{parskip}

\setlist{nosep,leftmargin=1.4em}

\title{\vspace{-2.5em}\textbf{Project Proposal}\\[0.3em]
\large Parallel Approximation of Betweenness Centrality with KADABRA}

\author{}
\date{}

\begin{document}

\maketitle
\vspace{-3em}

\begin{center}
\begin{tabular}{@{}ll@{}}
\textbf{Course} & Parallel and Distributed Computing (CS-3006), Fall 2026 \\
\textbf{Milestone} & 2, Project Proposal (due 23 November 2026) \\
\textbf{Base paper} & Paper 44: KADABRA is an ADaptive Algorithm for Betweenness \\
 & via Random Approximation (Borassi and Natale, ESA 2016) \\
\end{tabular}
\end{center}

\vspace{0.5em}

\begin{center}
\begin{tabular}{@{}llll@{}}
\toprule
\textbf{Role} & \textbf{Name} & \textbf{Roll number} & \textbf{Section} \\
\midrule
Contact person & Ibraheem Farooq & 23i-0816 & D \\
Member 2 & Musab Saleem & 23i-0857 & D \\
\bottomrule
\end{tabular}
\end{center}

\vspace{1em}

\section{Understanding the paper}

\subsection{The problem}

Some nodes in a network sit on the routes between many other pairs of nodes. Think of a bridge that connects two halves of a city. Betweenness centrality is the number that captures this idea. For a node $v$ it is defined as
\[
b(v) \;=\; \sum_{s \neq v \neq t} \frac{\sigma_{st}(v)}{\sigma_{st}},
\]
where $\sigma_{st}$ counts the shortest paths between $s$ and $t$, and $\sigma_{st}(v)$ counts how many of those paths pass through $v$. A node with a high score is a node that a lot of traffic has to flow through. This matters in social networks, road networks, and communication networks, where the high scoring nodes are the ones you want to protect, watch, or remove.

The trouble is the cost. The standard exact method, Brandes' algorithm, runs a shortest path search from every single node. That gives a running time proportional to the number of nodes times the number of edges. On a small graph this is fine. On a graph with millions of nodes it becomes far too slow to be useful, and there is strong theoretical evidence that nobody can do the exact computation much faster. This is why people turn to approximation.

\subsection{What KADABRA does}

KADABRA does not look at every pair of nodes. Instead it takes samples. It picks a random pair of nodes, finds a shortest path between them, and gives one point of credit to every node that sits on that path. It repeats this many times and divides by the number of samples. The result is an estimate of each node's betweenness. The more samples you take, the closer the estimate gets to the truth.

The clever part is knowing when to stop. KADABRA is adaptive, which means it does not fix the number of samples in advance. It works that number out as it goes, from the samples it has already seen. Here is how. For every node it keeps a shrinking error bar: given the samples so far, it works out how far above and how far below the true value the estimate could still be. These two bounds are usually written as the functions $f$ and $g$. The algorithm stops the moment every node's error bar is narrower than the target error, which we call $\epsilon$.

Two more pieces make this trustworthy. First, there is a ceiling on the number of samples, written $\omega$, fixed in advance from an upper bound on the graph's diameter. A graph whose nodes all sit close together needs fewer samples; one that is stretched out needs more. Second, the single allowed chance of being wrong, written $\delta$, is divided among the nodes, so each node is given a small share of it. A careful probability argument ties these together, so that stopping early, off the data, still delivers the promised result: with probability at least $1 - \delta$, every node's estimate lands within $\epsilon$ of the truth. An easy graph reaches that point quickly, and a hard one keeps sampling, up to $\omega$ if it must.

The paper also speeds up the sampling itself. Rather than searching the whole graph for each sample, it grows the search from both endpoints at once and stops when the two sides meet. This balanced two sided search finds a shortest path while touching far fewer nodes.

\subsection{A critical look}

KADABRA is a strong and widely cited method, but it is worth being clear about its limits. The accuracy guarantee is probabilistic, not absolute, so there is always a small chance a run is less accurate than promised. The stopping rule depends on an estimate of the graph's diameter computed beforehand, and a poor estimate costs either accuracy or time. The sampling is also randomized, so two runs give slightly different answers. For our project this is the honest picture we will keep in mind: the method is fast and principled, but it trades exactness for speed and relies on a careful stopping rule to stay trustworthy.

\section{Scope, challenge, and complexity}

\subsection{What we will build}

Our backbone is a later paper that parallelized KADABRA: van der Grinten, Angriman, and Meyerhenke (2019), \emph{Parallel Adaptive Sampling with almost no Synchronization}. It starts from the obvious parallel version, shows why the barrier in it limits scaling, and replaces that barrier with a lock free scheme based on epochs. Its results were measured on a server with 32 to 36 cores and large graphs. We do not have that machine. We have a laptop with a handful of cores and small graphs, and that gap is the thing we want to study.

So we treat the 2019 result as something to test, not something to repeat. On that large machine the lock free version clearly beats the barrier version. The honest question for our setting is whether that ordering still holds when there are only a few cores and the graphs are small, and if it does not, to say clearly where it breaks and why.

\textbf{Stage one, our committed deliverable.} We will write, in this order:
\begin{itemize}
  \item A sequential sampler, together with an exact Brandes computation to check it against. The sampler draws a random pair, samples one shortest path, and adds up the credits. We draw the path with a one sided search and weighted backtracking, described in Section 4.1.
  \item A straightforward OpenMP version with a barrier after each batch before the stopping check. This is the correct and simple baseline, and it is the version whose accuracy guarantee we can stand behind without any extra argument.
  \item One lock free version, the local frame scheme from the 2019 paper, where each thread keeps its own totals and publishes them in rounds, so the check never forces a barrier.
\end{itemize}
Both of those parallel versions are existing methods, the OpenMP baseline and the local frame scheme, both from the 2019 paper. We write them ourselves, which earns the marks for a correct parallel implementation, but writing them is not what we claim as new. They are the objects we measure. Our contribution is a careful characterization of how the barrier version and the local frame version behave in our regime. Across one, two, four, and eight threads, on small graphs, where does the lock free version help, where does it break even, and where does the simple barrier version actually win. The 2019 scaling numbers suggest much of the lock free advantage is a high core count effect, so a likely and honest finding is that the ordering flattens or reverses on a laptop. We will explain the causes we can measure: the cost of moving between rounds, the one serial step where a single thread adds up everyone's totals, and the fact that each sample on a small graph is so cheap there is little to overlap.

\textbf{Stage two, only if stage one is finished and solid.} The one thing the 2019 paper never varied is how its round length should depend on the graph itself. Its rule sets the round length from the thread count alone. We would sweep the round length and ask whether the best setting instead tracks the cost of a single sample, which grows with the graph's diameter. To keep this clean we would change the diameter on one synthetic graph generator while holding the node count, the edge count, and the degree fixed, so diameter is the only thing moving. This is the genuinely new question, and it is the stretch goal, not the floor.

\subsection{What we will leave out}

We are deliberately not doing several things. We will implement only one of the 2019 lock free variants, the local frame scheme, rather than the full set. We treat that re-implementation as a baseline to measure, not as our contribution, since the guide is clear that reproducing an existing method does not by itself count. We will not use the balanced two sided search that makes KADABRA's sampling sublinear, so we are clear that our speedups come from parallelism and not from fast path sampling. We will not use MPI or a cluster, because the message passing version from the 2020 follow up only pays off on many machines and large graphs, neither of which we have. We will not write a GPU version, because the irregular searches map poorly to a GPU and the effort is far more than two weeks. We will keep our graphs small enough that exact Brandes can check them, so we do not lean on any outside reference implementation. We will not try to beat the newest methods, such as SILVAN (2023), on sample count, since that lane is already well covered; our work is a systems and regime study, not a claim to estimate betweenness more accurately. Finally, we will not invent our own stopping rule, because a home made rule would break the accuracy guarantee and force us to prove a new one.

\subsection{Why this is challenging}

The difficulty here is not the mathematics, it is the parallel behaviour. Different samples take different amounts of time, so threads can end up waiting on each other. The shared table of scores is touched by every thread, which invites both race conditions and slowdowns from cache contention. Hardest of all, the stopping check has to look at a consistent set of totals. If it reads numbers that are halfway through an update, the guarantee quietly breaks and the output looks wrong for reasons that are hard to track down. This is more treacherous than it sounds, because our laptops use a processor with strong memory ordering, so a missing memory fence in the lock free version can pass every ordinary test and still be wrong. Only our accuracy check, repeated across many seeds, would catch it. Two smaller but real things also need care. The stopping rule needs an upper bound on the graph's diameter, and if that bound is ever too small the guarantee is void, so we compute it safely. And on a small graph the one time setup, computing that bound and a first calibration batch, can take longer than the sampling we are trying to measure, so we time the sampling phase on its own. Getting the lock free version correct is the real test of the project.

\subsection{Why the scope fits}

The sequential sampler and the barrier version are a few hundred lines each, which leaves our time for the harder parts: getting the lock free version correct, which earns the parallel implementation marks, and the measurement study, which is our actual contribution. Exact betweenness on a small graph is cheap to compute, so we always have a ground truth to check against. The course guide is clear that a careful, well explained result earns credit even when the speedup on small graphs is small or negative, so our plan is sized to produce honest measurements and a clear explanation rather than a headline number.

\section{Implementation plan and attribution}

\subsection{Parallelization strategy}

Threads draw independent samples in parallel. The shared state is the table of per node scores and a counter of samples taken. The one place that needs care is the stopping check, where the threads have to agree on a consistent set of totals before the rule is evaluated.

The barrier version reaches that agreement the simple way. Every thread finishes the batch, they meet at a barrier, and then one check runs on the settled totals. Because the barrier forces a clean, sequential looking state, this version's accuracy guarantee holds with no further argument, and we treat it as our reference for correctness.

The local frame version avoids the barrier. Each thread keeps its own copy of the totals and advances in numbered rounds, and a thread never writes to a copy that another thread is reading. When a check is due, the threads move to a new round and the check reads only the fully published previous round, using release and acquire ordering so the reader sees a complete set of writes. This is the subtle part, and it is where our care goes.

Two correctness rules hold in both versions. First, the stopping decision and the estimate we report use exactly the same set of samples, the ones in the snapshot the check read. We never fold in samples that were still being computed when the snapshot was taken. Second, that snapshot is always a whole number of finished samples taken at a round boundary, never a cut decided by the clock, because cheaper samples finish first and a snapshot taken by time would over represent short paths and bias every node's score.

For the path itself we run a one sided search from the first endpoint, recording for each node how many shortest paths reach it, then walk back from the second endpoint choosing each predecessor $p$ of the current node $u$ with probability $\sigma_{sp}/\sigma_{su}$. Those weights are what make the chosen path uniform among all shortest paths, and getting them wrong is the classic silent bug, so we test the sampler on a small graph with many equal length paths before trusting it.

\subsection{What we are contributing}

Our contribution is not to repeat the 2019 result but to test where it holds. Re-implementing the local frame scheme is prior work that we rebuild so we can measure it, and it is not the novel part; what is ours is the finding below. The paper showed, on a 32 to 36 core server, that the lock free version clearly beats the barrier version. We state a falsifiable claim for our own setting: at one, two, four, and eight cores on small graphs, that advantage largely disappears, and on the smallest and cheapest graphs the simple barrier version can match or beat the lock free one. We will report the point, in cores and in graph size, where the ordering changes, and we will tie the behaviour to specific causes we can measure, namely the latency of moving between rounds, the serial step where one thread sums everyone's totals, and the small cost of a single sample. Whether the result confirms or overturns our claim, it is a finding, because the course credits a well measured and well explained result even when the speedup is small or negative.

If stage one leaves time, the stretch question is sharper and genuinely open. The 2019 round length is tuned from the thread count alone. We would ask whether the best round length instead moves with the cost of a single sample. Our prediction, from the shape of the algorithm, is that a higher diameter makes each sample more expensive, which makes the check relatively cheaper, so the best round length should fall as diameter rises. We would test that by changing diameter on one synthetic generator with the node count, the edge count, and the degree held fixed, and by plotting the best round length against the measured average cost of a sample.

Throughout, we keep KADABRA's stopping rule unchanged, so the accuracy guarantee still holds and any error we report belongs to the sampling, not to a broken rule.

\subsection{Prior work we are comparing against}

The guide asks us to name any later paper that already implemented our approach and to treat it as our comparison baseline. Two papers apply:
\begin{itemize}
  \item van der Grinten, Angriman, and Meyerhenke (2019), the OpenMP version, which is our main backbone.
  \item van der Grinten and Meyerhenke (2020), the MPI version, which we mention for context but do not implement.
\end{itemize}
The 2019 paper reported roughly a sixfold speedup for the simple OpenMP version on 32 cores, and up to about three times more for its lock free variants. Those are large machine numbers. We will not copy any of them. Every figure in our report will come from our own runs on our own graphs, and the paper's numbers appear only as the published context our study sits beside. We also note SILVAN (2023), which needs fewer samples than KADABRA, to show we know where the field has moved; we are not competing with it, since our question is about parallel behaviour, not sample count.

\subsection{Existing code}

A reference implementation of KADABRA exists at \url{https://github.com/natema/kadabra}, released under the Apache 2.0 licence, and the 2019 algorithm ships inside NetworKit. We may read these for reference and will acknowledge them and keep their licence notices. The code we submit will be our own, and simply cloning or repackaging a repository is not our contribution.

\subsection{How we will evaluate}

\textbf{Correctness.} KADABRA is a randomized estimator, so a correct run does not match the exact answer, it stays within the promised error with high probability. We check it that way. Our ground truth is our own exact Brandes computation, first on a tiny graph we can read by hand. Then, over many seeds, we record the largest per node error against the exact scores, and the fraction of runs in which that error goes past the promised bound, and we check that this fraction stays at or below $\delta$, reporting it with a Clopper and Pearson confidence interval rather than a bare rate. We also report an average error and a ranking agreement, which tells us whether the important nodes come out in the right order. Before any comparison we put our estimates and the exact scores in the same normalization, dividing by $n(n-1)$ and fixing once whether pairs are counted in one direction or both, so a units mismatch cannot look like a real error. This same accuracy check is our only reliable guard against a missing fence in the lock free version. If that version ever fails it, we report its timings as informal and lean on the barrier version for any claim that rests on the guarantee.

\textbf{Performance.} We report speedup and efficiency,
\[
S = \frac{T_{\text{serial}}}{T_{\text{parallel}}}, \qquad E = \frac{S}{P},
\]
where $P$ is the number of threads, swept across one, two, four, and eight cores. We time the sampling phase on its own, apart from loading the graph, computing the diameter bound, and the first calibration batch, since on small graphs that setup can otherwise dominate. We use two timing modes, for two different questions. To compare the barrier and lock free versions cleanly, we run a fixed number of rounds so both do identical work, with the stopping check performed at every round boundary so its synchronization cost is really paid, but not used to end the run. That removes the randomness of adaptive stopping and leaves only the mechanism we are comparing. To see real behaviour, we also run with the adaptive rule switched on, and measure the wall clock time to reach the target accuracy together with how many samples were taken past the point where the run could have stopped. Because this second mode is random, we repeat it over many seeds and report the sample counts next to the times. We take at least twenty interleaved repetitions, report medians with their spread, run on mains power with the processor's boost disabled and threads pinned to physical cores on a machine whose cores are all the same kind, name the exact processor, and make no claim about larger core counts than the one we ran on.

\subsection{Risks and how we handle them}

\begin{itemize}
  \item \textbf{The speedup is small or vanishes at a few cores.} Expected, because the 2019 gain is largely a high core effect. We present a well measured null as the finding, and pick the thread count and graph where any effect is clearest.
  \item \textbf{The lock free version is subtly wrong on strong ordering hardware.} A missing fence can pass ordinary tests. The accuracy check over many seeds is our guard, and if it fails we fall back to the barrier version for any guaranteed claim.
  \item \textbf{Races on the shared table, or an inconsistent check.} We use a consistent snapshot and check the single threaded version against the exact answer before trusting the parallel one.
  \item \textbf{Biased path sampling.} We use the explicit $\sigma_{sp}/\sigma_{su}$ weights and test uniformity on a graph with many tied shortest paths.
  \item \textbf{A diameter bound that is too small voids the guarantee.} We compute a safe upper bound.
  \item \textbf{The stretch study confounds diameter with other structure.} We change diameter on one generator with size and degree fixed.
  \item \textbf{Random seeds across threads.} We fix a master seed and give each thread its own well separated stream.
  \item \textbf{We run out of time.} The barrier version is our insurance. On its own it satisfies the baseline and the parallel implementation marks, so the lock free version can slip without sinking the project.
\end{itemize}

\subsection{Who does what, and when}

We are two people with about two weeks. In the first week both of us build and validate the sequential sampler, the exact Brandes check, and the diameter bound, because everything else depends on them, and we add the one day OpenMP barrier version as soon as they pass. In the second week one of us attempts the local frame lock free version while the other builds the measurement harness and runs the experiments, and we gate the lock free version on the accuracy check before trusting any of its numbers. The stretch round length study is attempted only once the barrier and lock free comparison is written up. Both of us will be able to explain every part of the code and the results.

\begin{thebibliography}{9}

\bibitem{kadabra} M. Borassi and E. Natale. KADABRA is an ADaptive Algorithm for Betweenness via Random Approximation. \emph{European Symposium on Algorithms (ESA)}, 2016. \url{https://drops.dagstuhl.de/entities/document/10.4230/LIPIcs.ESA.2016.20}

\bibitem{vdg2019} A. van der Grinten, E. Angriman, and H. Meyerhenke. Parallel Adaptive Sampling with almost no Synchronization. \emph{Euro-Par}, 2019. \url{https://arxiv.org/abs/1903.09422}

\bibitem{vdg2020} A. van der Grinten and H. Meyerhenke. Scaling Betweenness Approximation to Billions of Edges by MPI-based Adaptive Sampling. \emph{IPDPS}, 2020. \url{https://arxiv.org/abs/1910.11039}

\bibitem{brandes} U. Brandes. A Faster Algorithm for Betweenness Centrality. \emph{Journal of Mathematical Sociology}, 2001. \url{https://doi.org/10.1080/0022250X.2001.9990249}

\bibitem{silvan} L. Pellegrina and F. Vandin. SILVAN: Estimating Betweenness Centralities with Progressive Sampling and Non-uniform Rademacher Bounds. \emph{ACM Transactions on Knowledge Discovery from Data}, 2023. \url{https://arxiv.org/abs/2106.03462}

\bibitem{repo} Reference implementation of KADABRA. \url{https://github.com/natema/kadabra}

\end{thebibliography}

\vspace{1em}
\noindent\footnotesize{AI assistance disclosure: AI tools were used to help research related work and draft and structure this proposal. All technical decisions and claims are the group's own, and every member is responsible for the submitted content.}

\end{document}
